\subsubsection{The dual-write replay}
\label{sec:replay}

Everything to this point compares a PromQL candidate against our reading of Datadog's
rollup. That is a model, and a model is what \S\ref{sec:not-measured} warns the reader
about. This subsection closes the gap for the counter rules by running the dual-write
construction of this section literally: one stream of observations, written to both
systems, read back from both.

The stream is real. Counting the requests in a web-server access log into a scrape grid
gives a counter with the shape that \S\ref{sec:threats} says the synthetic workloads only
imitate: in the half hour used here the per-minute rate swings seventeenfold without
being asked to. We take the busiest half hour of the Kennedy Space Center HTTP logs of
July~1995~\citep{nasahttp}, count requests at a $15$\,s scrape interval into one counter
per HTTP status class, and write the cumulative series to Prometheus and its image under
$\phi_{\mathsf{zero}}$ --- the per-interval deltas --- to Datadog, at the same timestamps
(\texttt{scripts\slash replay\_dualwrite.py}).
The trace's publishers ask that it not be analysed beyond general traffic patterns, and
the parser enforces that rather than promising it: it keeps a timestamp and a status
class and discards the host and the request line, so what is replayed is a count of
requests per interval and nothing else. Timestamps are shifted to the present, because
Datadog's intake refuses points much older than an hour; the shape is real and the
absolute times are not.

\begin{table}[htbp]\centering\footnotesize
\begin{tabular}{@{}llrrr@{}}
\toprule
Scrape grid & PromQL candidate & Evaluated inside & On the edge & Max $d$ \\
\midrule
unaligned & \texttt{x - x offset $\Delta$} & 100.0\% & 100.0\% & 0 \\
 & \texttt{increase(x[$\Delta$] anchored)} & 100.0\% & 100.0\% & 0 \\
 & \texttt{rate(x[$\Delta$])} $\times\,\Delta$ & 20.2\% & 20.2\% & 1 \\
\addlinespace
aligned & \texttt{x - x offset $\Delta$} & 100.0\% & 27.4\% & 0 \\
 & \texttt{increase(x[$\Delta$] anchored)} & 100.0\% & 27.4\% & 0 \\
 & \texttt{rate(x[$\Delta$])} $\times\,\Delta$ & 20.2\% & 16.7\% & 1 \\
\addlinespace
\bottomrule
\end{tabular}

\caption{The dual-write replay: each candidate against \emph{Datadog}, bucket by bucket,
on a real request trace. ``Inside'' and ``on the edge'' are the two evaluation instants
of Remark~\ref{rem:closure}. $84$ comparisons per row.}
\label{tab:replay}
\end{table}

Table~\ref{tab:replay} is the result, and two things in it are worth separating.

The first is that both exact forms are exact --- against Datadog, not against a model of
it, on traffic nobody generated for the purpose. \texttt{rate} is not, at $20.2\%$ of
buckets on both grids, which is what \S\ref{sec:ingestion} has said throughout and what
the steady synthetic workload conceals by being steady. That figure used to read
$19$--$20\%$, a spread between the two grids that was not a property of either: the
replay script scored this one table by bit-identity where every other exactness column in
this paper uses a $10^{-12}$ relative bound, and one aligned bucket reading $236.0$
against $235.999\ldots$ counted as a disagreement.

The second is about the convention of Remark~\ref{rem:closure}, and it explains how that
convention came to be assumed wrongly for so long. On the unaligned grid --- a scrape
whose samples do not fall on rollup boundaries, which is what a jittered scrape produces
and therefore the ordinary case --- evaluating on the bucket's edge agrees with
evaluating inside it, on every bucket. Off a boundary the two half-open intervals contain
the same samples, so the closure is simply not observable. Snap the scrape grid onto the
rollup grid and the same comparison falls from $100\%$ to $27.4\%$. The error was
invisible in exactly the configuration real deployments have, and decisive in the one the
synthetic workloads of Table~\ref{tab:rollup-live} happen to use.

What this does not do is also worth stating. A request counter does not reset in half an
hour, so \texttt{COUNTER\_RESET} is not exercised here; the restart rows of
Table~\ref{tab:rollup-live} and the scripted reset of E1 (Remark~\ref{rem:rho}) are the
evidence for that case. It is one window of one
trace at one rollup interval. And it tests the counter rules, not translation coverage:
$\kappa$ is what it was.

\paragraph{What was assumed, and what is left.}
This section began with two facts about Datadog that its documentation does not settle,
and carried a third about its backend. None is still an assumption, and the two it began
with did not turn out the way they were guessed.

The reset policy $\rho$ is no longer among them either, and it is the one that was
settled twice: read from the agent's source and then put to a running agent, which
agreed with the reading (Remark~\ref{rem:rho}). It too went against the guess --- the
prediction recorded in advance was discarding the interval, which turns out not to be a
reset policy at all.

Bucket closure is no longer among them. Remark~\ref{rem:closure} settles it by
measurement: buckets are left-closed and named by their start, which is the opposite of
what an earlier draft of this section assumed and the opposite of PromQL's own
convention. Epoch alignment \emph{is} documented~\citep{ddrollup}, so the remaining
content of \texttt{ROLLUP\_ALIGNMENT} is sharper than ``align the grids'': a PromQL
window follows the evaluation step, and the step that reads a Datadog bucket is one
\emph{inside} its right edge rather than on it. Evaluating on the boundary --- the
obvious choice once the closure is known --- disagrees on every bucket where a scrape
lands on the boundary too, which is the ordinary case. The translator records
\texttt{ROLLUP\_ALIGNMENT} on every rollup it emits for that reason.

Nothing is left. The last of the three --- whether Datadog's backend stores what its
agent sends --- stood until a counter was made to reset on purpose, which is what the
replay above cannot do and what E1 does (Remark~\ref{rem:rho}). Both systems scrape one
endpoint across the reset; Datadog stores $103$ over the bucket holding it, and
\texttt{increase($x$[$60$s] anchored)} over Prometheus returns $103$ for the same bucket,
as it does on every complete bucket of all three counters. So the restart rows of
Table~\ref{tab:rollup-live} no longer stand alone: the anchored form tracks
$\phi_{\mathsf{zero}}$, and Datadog agrees with it where a counter actually restarts.

\iffull\apphead{What the vendor was asked}{app:experiments}

Datadog does not document the behaviour several of this paper's claims rest on, so it was
measured against a trial account and the raw responses cached under \texttt{data/raw/}.
The last of these asks the vendor's published integration repository rather than its API,
because what the dashboards omit that repository declares.
Each experiment below states its design, its prediction table where it has one, and what
came back; the body cites the answers.

\begin{remark}[Which interval a bucket covers, measured]\label{rem:closure}
That the bucket is \emph{left}-closed and named by its \emph{start} is measured rather
than assumed, and it is not what we first assumed. Submitting one point at an exact
multiple of $\Delta$ and one strictly inside a bucket, and reading back which bucket each
is reported under, separates the four conventions: a point on a boundary is reported
under the bucket \emph{starting} there. Each of the other three conventions is ruled out
by at least one of the two readings (\texttt{scripts\slash trial\_checks.py e2}; the
responses are cached and regression-tested). It is one run at one rollup interval ---
the script fixes $\Delta$ at $60$\,s --- so what is pinned is the convention at that
width, not its independence of the width.

The convention PromQL uses for its own range selectors is the opposite one, $(u-w,u]$,
which is why it was the natural guess. Getting it wrong cost nothing numerically --- a
model and a candidate evaluated under the same convention agree with each other, so every
figure in Table~\ref{tab:rollup-live} is what it was --- but it was a claim about Datadog
that happened not to be true, and it changes where a translated query must be evaluated.
\end{remark}

\begin{remark}[Which $\rho$, and how we know]\label{rem:rho}
Datadog's documentation does not say what becomes of a counter decrease. Its agent is
open source, so the question can be answered by reading it~\citep{ddagent}, and then put
to a running one; both were done and they agree. The reading first.

The conversion carries a per-sample flag. With the flag set, a
sample below its predecessor contributes its own value, on the agent's stated assumption
that the counter restarted from zero --- this is $\delta_{\mathsf{zero}}$ term for term.
With the flag clear, the decrease contributes nothing and the point is still emitted,
which is $\delta_{\mathsf{clamp}}$. The OpenMetrics check sets the flag after every
successful scrape, so the ingestion path this paper models is $\phi_{\mathsf{zero}}$ from
its second scrape onward: the policy PromQL's own reset handling assumes.

Both branches above are thus implemented, and a third that an earlier draft of this
section offered --- discarding the interval outright --- is not, at least not as a reset
policy. Discarding is instead what happens at the \emph{start} of a series, where there is
no predecessor to difference against. A failed scrape is the interesting case: it clears
the flag, so the next interval is converted under $\mathsf{clamp}$ rather than
$\mathsf{zero}$, and is dropped altogether if the aggregator's stored predecessor has
expired in the meantime. Either way a transient failure perturbs $\phi(D)$ in a way
nothing applied to $D$ reproduces, and Proposition~\ref{prop:admissible} gives no warning,
because the discrepancy is between the two stored databases rather than between two
queries. No refusal can flag it and no translation can repair it.

Reading the agent establishes what it intends to send, not what the backend stores, and
the two are different claims. E1 closes the gap by putting a real agent in the loop. A
local endpoint exposes a counter that climbs by one per scrape to $1000$ and then, inside
one rollup bucket, drops to $100$; a Datadog Agent runs an OpenMetrics check against it
and a Prometheus scrapes the same endpoint. Over the bucket holding the reset Datadog
stores $\mathbf{103}$ across four points --- the increment before the reset, the new
value reported as the reset interval's own count, and the two increments after it. That
is $\delta_{\mathsf{zero}}$ and nothing else: $\mathsf{clamp}$ predicts $3$ over four
points, discarding the interval predicts $3$ over three, and a backend keeping the signed
difference predicts $-897$. The four predictions were recorded before the run and are
regression-tested in \texttt{tests\slash test\_reset.py}. A second counter resetting to
exactly $0$ rather than $100$ reads $3$, which is what $\mathsf{zero}$ requires and rules
out a special case at zero; a third that never resets reads four increments of one in
every bucket, which is what shows the run is reading the bucket it means to.
\end{remark}

\begin{remark}[$\phi_\rho$ and the flush interval]\label{rem:flush}
$\phi_\rho$ differences consecutive \emph{samples}, and the agent aggregates samples
into the interval it flushes. An earlier version of this remark inferred from the source
that the two part company as soon as a flush holds more than one sample: a decrease
appeared to \emph{overwrite} what the interval had accumulated rather than adding to it,
so a flush holding an increase and then a reset should report the post-reset value alone,
where $\phi_{\mathsf{zero}}$ keeps both. The agent's default scrape and flush intervals
are both $15$\,s, which would make the question moot in the default configuration and
was the reason it went untested.

Put to a running agent, it does not happen --- though not quite for the reason the run
was designed to test, and the difference is worth keeping. Driving the OpenMetrics
check at a $5$\,s collection interval against the unchanged $15$\,s flush, the bucket
holding the reset reports $111$: eleven increments of one and the reset's own value,
which is $\phi_{\mathsf{zero}}$ term for term. Losing the increment that shares a flush
with the reset would give $110$ (\texttt{scripts\slash reset\_policy.py -{}-step 5}).

What the stored data then shows is that no flush combined anything. Every complete
$60$\,s bucket holds twelve points, one per collection, where three samples to a flush
would have left four. So the agent forwarded each sample rather than aggregating within
the flush, and the overwrite this remark used to posit had no opportunity to occur. The
hypothesis is therefore unexercised rather than refuted, and what the run does
establish is the end-to-end fact the exactness claims actually use: at this ratio
Datadog stores $\phi_{\mathsf{zero}}$ term for term, so those claims do not depend on
the scrape and flush intervals agreeing. The ratio itself was configured and not
observed --- the record computes it from the two interval settings --- and no run here
drives it to a point where a flush must combine samples.
\end{remark}

\begin{remark}[\texttt{anchored} is still experimental]\label{rem:anchored-flag}
The \texttt{anchored} modifier of~\citet{prom0052} arrived in Prometheus~3.7.0 behind
\texttt{-{}-enable-feature=promql-extended-range-selectors}, and an earlier draft of this
section asserted that the flag had since become a no-op. It has not. A $3.15.0$ server
started without it rejects the expression outright:
\begin{center}\footnotesize
\texttt{anchored modifier is experimental and not enabled}
\end{center}
The precondition \texttt{ANCHORED\_VERSION} records is therefore stronger than a version
number. It is a server deliberately started with an experimental feature enabled, which
is a different and larger ask of whoever operates the target: a version can be waited
for, a flag has to be chosen. That is the real cost of
Proposition~\ref{prop:anchored}'s exactness, and it is the reason the emitter leaves the
reset-aware form off by default and emits the cumulative difference, which uses nothing
outside standard PromQL.

The error above is how we found out. Every measurement in this paper had been taken with
the flag set, because the script that takes them passes it and a comment there described
it as harmless; the claim was never exercised until a later experiment started a server
without it. A precondition that every instrument satisfies by habit is one nobody tests.
\end{remark}

\begin{remark}[Where \texttt{.as\_rate()}'s divisor sits]\label{rem:as-rate-placement}
The rollup is inside the space aggregation and the modifier's division is not, which is
a second ordering question and was answered the same way round as the first: wrongly, at
first, and then by measurement.

\texttt{.as\_rate()} divides by the rollup interval. Dividing by a positive constant
commutes with $\mathsf{sum}$, $\mathsf{avg}$, $\mathsf{min}$ and $\mathsf{max}$, so for
those the divisor can sit on either side of the space aggregation and nothing
distinguishes the readings. It does not commute with $\mathsf{count}$, which counts
series and is blind to their values: $\mathrm{count}(x/\Delta) = \mathrm{count}(x)$
while $\mathrm{count}(x)/\Delta$ is not. So one operator decides the question for all
of them, and asking it of a real account answers it ---
\texttt{count:m\{*\} by \{g\}.rollup(avg, 30)} returns $1$ per group and the same
query with \texttt{.as\_rate()} returns $1/30$
(\texttt{scripts\slash as\_rate\_order.py}). Datadog divides the aggregated result, so
the emitter places the divisor outside the aggregation for $\mathsf{count}$ and, for
readability, leaves it inside where it commutes.

The same run settles the modifier's other ordering, against an explicit
\texttt{.rollup}. Submitting every $10$\,s and rolling up over $30$\,s separates
dividing each delta by its own submission interval from rolling up and dividing once by
the bucket: the first reads $1$, the second $10/30$. Datadog reads $10/30$. Both
readings coincide whenever the submission grid equals the rollup width, which is every
harness in this project and the reason the question went unasked for so long.
\end{remark}

\begin{remark}[The assumption that underwrites all of the output]%
\label{rem:gauge}
Writing Theorem~\ref{thm:kernel} down is what exposed this. An
$\mathit{op}\in\{\mathsf{avg},\mathsf{min},\mathsf{max},\mathsf{count}\}$ rollup reduces
the points Datadog \emph{stores}. For a GAUGE those are the samples Prometheus also
holds, and \texttt{avg\_over\_time(x[$\Delta$s])} is exact. For a COUNT they are
per-interval deltas, and averaging deltas is not averaging running totals --- no
\texttt{*\_over\_time} of the cumulative series computes it, and the emitted expression is
not approximately right but unrelated. The dashboard never says which, so the emission is
correct for a gauge and the emitter records \texttt{METRIC\_TYPE\_UNKNOWN} to say so. It
reports rather than refuses
because a rule \emph{can} be chosen here; for a $\mathsf{sum}$ rollup none can, and the
same code is a refusal there instead.

The condition is not a precaution. Asked of a real account, a bare query on a
\texttt{COUNT} metric returns the \emph{mean} of four stored deltas,
so Datadog applies its enforced \texttt{avg} default rather than the \texttt{.as\_count()}
its user interface appends --- averaging deltas where \texttt{avg\_over\_time} averages
running totals (\texttt{scripts\slash trial\_checks.py e3}). The deltas are $1$, $14$,
$2$ and $3$, chosen so that each candidate predicts a different reading: $\mathsf{avg}$
gives $5$, $\mathsf{sum}$ gives $20$, a $\mathsf{count}$ rollup gives $4$ and
$\mathsf{sum}$ over the reported interval gives $1$. The account returns $5$. The first
version of this probe used $1$, $2$, $3$ and $10$, whose mean is $4$ and whose count is
also $4$, and named $\mathsf{avg}$ against a rival set that did not include
$\mathsf{count}$ --- so the reading it took could not have told the two apart.

The condition was also silent until now, and it is not rare. At an earlier stage of this
work it appeared to hold of $6{,}381$ translated occurrences --- $76.2\%$ of the
translated set as it then stood, which was smaller than today's --- because Datadog's
default time aggregator is $\mathsf{avg}$ and most queries state no method at all. The remaining quarter looked exempt: a $\mathsf{sum}$ rollup needs a
counter, the ways to ask for one are \texttt{.as\_count()}, \texttt{.as\_rate()} and
\texttt{.rollup(sum, $\Delta$)}, and all three are documented for \texttt{COUNT}
metrics --- so a query using one, we reasoned, establishes the type it does not state.

It does not, and the error was of a kind this paper has made before: a rule read in a
manual and never asked of the system. Submitting a \texttt{GAUGE} and querying it with
each modifier returns the plain rollup, unchanged and undivided --- Datadog accepts both
on a gauge and silently ignores them rather than rejecting the query
(\texttt{scripts\slash as\_rate\_order.py}). The modifier therefore records what the
dashboard's author took the metric to be. It is evidence about a person, not about a
type, and no reading of the query distinguishes the two.

The third way is weaker still, and for a different reason. \texttt{.rollup(sum, $\Delta$)}
over a \texttt{GAUGE} is not ignored: the same experiment reads it as the sum of the
stored samples, which is a perfectly meaningful query that a dashboard author might well
intend. So where \texttt{.as\_count()} on a gauge is a mistake nobody benefits from,
summing a gauge over a window is an ordinary thing to want --- and the translation this
front end chooses for it, the cumulative difference, computes something else entirely.
Of the three pieces of evidence the query can carry, the one that is hardest to dismiss
as author error is the one whose misreading is least benign. With that premise withdrawn the
caveat holds of all $8{,}475$ translated occurrences, $100\%$ of them
(Table~\ref{tab:translation-caveats}).

Recording it does not move $\kappa$ --- the translations are unchanged --- which is the
point: it changes what they are claimed to be. Its refusing counterpart still never fires,
but for a weaker reason than the one we gave: not because the method establishes the type,
but because it is the only syntax that asks for the rule, so the emitter always has a rule
to choose and never has to stop. So the same unstated fact is a diagnostic that cannot
occur and a precondition that always does.
\end{remark}

\begin{remark}[What hypothesis~(3) costs when it fails]\label{rem:type-cost}
Hypothesis (3) is the one that changed shape while this paper was being written, and
twice. It began as ``the metric is a gauge'', which was the only case the emitter then
handled; a counter under an averaging rollup was being emitted as
\texttt{avg\_over\_time} of the running totals, which is not a function of the deltas
Datadog reduces at all, and the corpus-scale differential test of \S\ref{sec:differential}
is what caught it. It then stopped being dischargeable from the query text, when the
modifiers it rests on turned out not to establish the type (Remark~\ref{rem:gauge}).
What it costs when it fails is measured rather than asserted: writing one stream of
counter observations to both systems and asking the commonest shape in the corpus ---
a query with no method at all --- Datadog reads the mean of the deltas, $10$ throughout,
while \texttt{avg\_over\_time} reads the mean of the running totals, climbing from $20$
to $1{,}760$ over half an hour. The relative difference runs from $0.50$ to $0.994$ and
is unbounded in the counter's age, because one quantity grows with it and the other does
not. The same run with the metric submitted as a \texttt{GAUGE}, where the hypothesis
holds, agrees exactly on every bucket (\texttt{scripts\slash type\_assumption.py}).
\end{remark}

\begin{remark}[How often the inference is wrong]\label{rem:type-measured}
A precondition that holds of all the output and cannot be checked is the weakest thing
this paper carries, and it turns out to be checkable. The dashboards do not record a
metric's type --- but the repository they come from does. Beside each integration's
\texttt{assets\slash dashboards\slash}, \texttt{integrations-core} ships a
\texttt{metadata.csv} with a \texttt{metric\_type} column, and
\texttt{scripts\slash fetch\_corpus.py} had taken only the dashboards.
\texttt{scripts\slash fetch\_metric\_types.py} fetches the rest, reading the commit
from the corpus's own record so the declarations cannot come from a different tree than
the queries. At that commit those files declare $37{,}070$ metric names: $23{,}597$
\texttt{gauge}, $13{,}140$ \texttt{count}, $330$ \texttt{rate} and three that two
integrations declare incompatibly.

Joining the two settles the hypothesis for most of the corpus
(\texttt{scripts\slash type\_declared.py}). Of the $8{,}444$ translated occurrences,
$7{,}557$ --- $89.5\%$ --- name a metric the repository declares as a gauge or a count,
and the front end's inference disagrees with the declaration on $491$ of them: $6.5\%$
of the decidable set, $5.6\%$ of the corpus. The two directions of error are both
present and are different failures. In $242$ occurrences the metric is a
\texttt{count} and the query carries no evidence of it, so the emitter translates it as
a gauge and the \texttt{*\_over\_time} it produces reduces running totals where
Datadog reduces deltas --- the case this remark describes as unrelated rather than
approximate. In $249$ the metric is a \texttt{gauge} and the query carries
\texttt{.as\_count()}, \texttt{.as\_rate()} or \texttt{.rollup(sum, $\Delta$)}
anyway, so the evidence is not merely weak but actively wrong, and the emitter chooses
a counter rule for a series that has no deltas to difference.

What this is evidence of, and what it is not. \texttt{metadata.csv} states what an
integration submits, so a disagreement says the inference is wrong, not that it must
be: a dashboard can query a metric the integration does not ship, and an operator can
submit the same name as something else. That is the same standard \S\ref{sec:emission}
applies to the modifier evidence this replaces, and it is a sharper instrument than the
one it replaces rather than a proof. The remaining $887$ occurrences divide three ways: $736$ name metrics the
repository does not declare, $134$ name \texttt{rate} metrics --- a third
submission type Definition~\ref{def:phi} maps by the identity rather than modelling ---
and $17$ name one of the three metrics two integrations declare incompatibly, so the
repository states a contradiction rather than a type.
\end{remark}

\fi
